\chapter{Healing: Mending Flesh and Spirit}
\label{ch:healing}
\index{healing}\index{Fatigue}\index{Harm}

\epigraph{"The Mend tag closes small wounds and eases exhaustion. To knit broken bones or purge a curse, you need greater power---and greater risk."}{---The Gray Wanderer}

In \textbf{Fate's Edge}, injuries are not mere hit points. They are narrative weights: a fractured leg forces hard choices; a festering wound breeds desperation. Healing magic exists, but it is never cheap, never perfectly safe, and always part of the story.

This chapter explains how your character recovers from the two kinds of damage---\textbf{Fatigue} and \textbf{Harm}---using magic, medicine, rest, and willpower.

\section{The Two Faces of Injury}
\label{sec:two-faces-injury}

\begin{table}[h]
\centering
\begin{tabular}{lll}
\toprule
\textbf{Type} & \textbf{What It Represents} & \textbf{How It Heals (Fastest)} \\
\midrule
\textbf{Fatigue} & Exhaustion, minor cuts, bruises, winded & Mend cantrip, short rest, a hot meal \\
\textbf{Harm} & Broken bones, deep wounds, poisoning, curses & Rest + medicine, powerful magic, time \\
\bottomrule
\end{tabular}
\end{table}

\paragraph{Fatigue} is your character's ability to keep going. It accumulates from physical exertion, missed sleep, harsh weather, and the strain of magic. Fatigue is \textbf{common and easy to remove}.

\paragraph{Harm} is serious injury. Each level of Harm (1, 2, or 3) imposes a −1 die penalty on most actions. Harm is \textbf{rare and expensive to heal}. When you take Harm, the world takes notice.

\section{Recovering Fatigue -- The Easy Work}
\label{sec:recovering-fatigue}

Fatigue clears quickly with rest, food, or a minor magical touch.

\subsection{Rest and Recovery}
\label{subsec:rest-fatigue}

\begin{itemize}
  \item \textbf{Short rest (one hour):} Remove 2 Fatigue. You must be able to sit, breathe, and drink water. No roll required.
  \item \textbf{Full night's rest (8 hours in safety):} Remove all Fatigue. If you are in a dangerous or exposed place, the GM may reduce this to 2 Fatigue instead.
  \item \textbf{Warm meal and comfort:} The first rest after a hot meal removes an additional 1 Fatigue (stacks with the above).
\end{itemize}

\subsection{The Mend Cantrip}
\label{subsec:mend-cantrip}

The \texttt{[HEAL]} tag (typically from the Life element or a Patron of healing) removes Fatigue quickly.

\begin{tcolorbox}[colback=green!5, colframe=green!75!black, title=Mend Cantrip]
\begin{tabular}{ll}
\textbf{Tags:} & \texttt{[HEAL]} \\
\textbf{DV:} & 2 \\
\textbf{Action:} & 1 action \\
\textbf{Range:} & Touch \\
\textbf{Effect:} & Remove \textbf{1 Fatigue} from yourself or a willing creature you touch. \\
\textbf{Backlash:} & None. The GM may still spend Story Beats (from any 1s on the roll) to add minor flavour (e.g., a faint scar, a shiver down the spine).
\end{tabular}
\end{tcolorbox}

\noindent This cantrip works for any magical tradition that can access Life or a healing Patron (e.g., Free Caster, Runekeeper of Oath of Flame, Cantor with healing songs).

\paragraph{Example:} Lyra uses Mend on Kael after a fight. She rolls 5d10 (Spirit + Medicine), gets 3 successes vs DV 2 -- a Clean Success. Kael clears 1 Fatigue. No complication.

\section{Healing Harm -- The Hard Work}
\label{sec:healing-harm}

Restoring a level of Harm (e.g., from Harm 2 to 1, or Harm 1 to 0) requires \textbf{significant effort} -- either mundane medical care over time or powerful magic that comes with risk.

\subsection{Mundane Healing}
\label{subsec:mundane-healing}

\begin{itemize}
  \item \textbf{First Aid (Wits + Medicine, DV 3):} Stabilise a dying creature (prevent Harm from worsening). Does not heal existing Harm. Takes 1 action.
  \item \textbf{Treatment (downtime, safe haven):} Requires a healer's kit and at least one day of rest. Restores 1 Harm per full day of care. A successful Medicine roll (DV 3) can reduce required days by half (minimum 1 day).
  \item \textbf{Surgery or magical cleansing (serious injury):} For Harm 3 or cursed wounds, the GM may require a full scene or a quest component (rare herbs, a spirit's favor).
\end{itemize}

\subsection{Magical Healing of Harm}
\label{subsec:magic-healing-harm}

To heal Harm with a spell, you must meet \textbf{two conditions}:

\begin{enumerate}
  \item Include the \texttt{[HEAL]} tag.
  \item Include \textbf{at least one other appropriate tag} based on the injury:
    \begin{itemize}
      \item \texttt{[PURIFY]} -- poison, disease, corruption
      \item \texttt{[STRENGTHEN]} -- broken bones, torn muscles
      \item \texttt{[CLEANSE]} -- curses, magical wounds, shadow‑damage
    \end{itemize}
\end{enumerate}

The base DV for healing Harm is \textbf{1 + number of tags}. Healing Harm also adds a \textbf{+2 DV surcharge}. Examples:

| Injury | Tags | Base DV | +2 Surcharge | Final DV |
|--------|------|---------|--------------|----------|
| Shallow cut (Harm 1) | \texttt{[HEAL]} | 2 | +2 | 4 |
| Poisoned wound (Harm 1) | \texttt{[HEAL][PURIFY]} | 3 | +2 | 5 |
| Broken arm (Harm 2) | \texttt{[HEAL][STRENGTHEN]} | 3 | +2 | 5 |

The GM may increase DV further for legendary wounds (e.g., a dragon's claw tears your spirit -- \texttt{[HEAL][CLEANSE]} + a ritual requiring significant time).

\subsection{Ritual Healing}
\label{subsec:ritual-healing}

Healing Harm 2 or Harm 3 often requires a \textbf{ritual}, not a single action. Rituals take at least 10 minutes (or longer per the GM). During a ritual, interruptions (combat, loud noises, disbelief) can cause the healing to fail or backfire.

Many magical paths have specific healing rites:
\begin{itemize}
  \item \textbf{Runekeeper (Oath of Flame \& Light):} \textit{Lay on Hands} (Low Rite) heals Fatigue; \textit{Cleansing Light} (Standard Rite) can remove poison or downgrade Harm at the cost of Obligation.
  \item \textbf{Cantor:} Healing songs (e.g., \textit{Hymn of Restoration}) require a Performance roll. On a Partial, the wound heals but the Cantor marks Corruption.
  \item \textbf{Witch (Hearth craft):} Hedge Gifts like \textit{Steady Hand} (remove Fatigue) and \textit{Knot of Mercy} (remove a Condition) are roll‑free but limited per scene/session.
  \item \textbf{Invoker:} Ritual healing with a Patron's Symbol is slow but has reduced backlash. Crack the Seal for emergency healing -- at double Obligation cost.
  \item \textbf{Summoner:} Can bargain with a spirit to transfer wounds (e.g., a spirit of sacrifice takes your Harm). Tick the \textbf{Leash} clock each time.
\end{itemize}

\section{The Risk of Magic -- Story Beats and Backlash}
\label{sec:magic-healing-risk}

Whenever you cast a healing spell (especially to heal Harm), any \texttt{1}s you roll generate Story Beats (SB) for the GM. The GM will spend those SB to introduce \textbf{complications} -- the healing works, but the world pushes back.

\paragraph{What complications look like:}
\begin{itemize}
  \item \textbf{Minor (1 SB):} The wound leaves a visible scar. The patient feels an odd craving (honey, raw meat, silence). A faint magical residue lingers.
  \item \textbf{Moderate (2 SB):} A minor spirit or "echo" of the injury appears -- it watches the patient for a scene. A Patron sends a dream: "You owe me."
  \item \textbf{Serious (3 SB):} The healing draws attention from a rival, a witch‑hunter, or a territorial nature spirit. The patient gains a thematic Condition (e.g., \textit{Phantom Pain}, --1 die on rolls involving that limb until they perform a small ritual).
  \item \textbf{Major (4+ SB):} Reality warps briefly -- a corpse nearby twitches; the patient forgets something precious; the healer's own shadow moves on its own for a day.
\end{itemize}

These complications \textbf{do not cancel the healing}. They are \textbf{narrative gifts} -- hooks for future scenes, reminders that power has a price.

\begin{tcolorbox}[colback=black!5, colframe=black!50, title=GM's Role]
Your GM uses Story Beats to make healing \textbf{dramatic}, not to punish you. Embrace the complications. A scar is a story. A Patron's demand is an adventure. A haunted shadow is a future ally (or enemy). Talk to your GM about the kinds of complications you find interesting -- body horror, social fallout, ominous omens, or debts to mysterious powers.
\end{tcolorbox}

\section{Natural Recovery -- When Magic Fails}
\label{sec:natural-recovery}

If you have no access to magic or medicine, your body still heals -- slowly.

\begin{itemize}
  \item \textbf{Per day of complete rest (no strenuous activity):} Remove 1 Fatigue automatically. If you also have food, water, and shelter, remove 2 Fatigue.
  \item \textbf{Per week of rest (downtime):} Remove 1 level of Harm. A successful Medicine roll (Wits + Medicine, DV 3) at the start of the week reduces the required time to 3 days per Harm level.
  \item \textbf{Neglected injuries:} If you continue to adventure while at Harm 2 or Harm 3 without any healing, the GM may advance a \textbf{Worsening Condition} clock. When it fills, your injury becomes permanent (a limp, a weakened arm, etc.) or requires extraordinary magic to fix.
\end{itemize}

\section{Healing by Magical Tradition -- Quick Reference}
\label{sec:healing-by-tradition}

\begin{table}[h]
\centering
\begin{tabular}{lll}
\toprule
\textbf{Path} & \textbf{Fatigue Healing} & \textbf{Harm Healing} \\
\midrule
Free Caster & \texttt{[HEAL]} cantrip (DV 2) & \texttt{[HEAL]}+other tag, +2 DV, risk backlash \\
Runekeeper & Low Rite (e.g., \textit{Lay on Hands}) & Standard Rite, costs Obligation, may Push \\
Invoker & Ritual \texttt{[HEAL]} (slow) & Ritual + Symbol, reduce DV with rare components \\
Cantor & Healing Song (Performance roll) & Song + Push (marks Corruption), requires High Cantor for instant \\
Witch (Hearth) & Hedge Gift: \textit{Steady Hand} (no roll, 1/scene) & Hedge Gift: \textit{Knot of Mercy} (once/session, removes Condition) \\
Summoner & Spirit bargain (costs Leash ticks) & Transfer wound to spirit (costs Leash, may anger spirit) \\
No magic & Short/long rest & Downtime + Medicine roll \\
\bottomrule
\end{tabular}
\end{table}

\section{Summary: Healing at a Glance}
\label{sec:healing-summary}

\begin{table}[h]
\centering
\begin{tabularx}{\textwidth}{lX}
\toprule
\textbf{What You Want} & \textbf{How to Get It} \\
\midrule
Clear 1 Fatigue (fast) & Mend cantrip (DV 2, 1 action) or short rest (1 hour). \\
Clear all Fatigue & Full night's rest in safety. \\
Heal 1 Harm (mundane) & A week of rest, or 1 day with Medicine roll (DV 3) and a healer's kit. \\
Heal 1 Harm (magic) & \texttt{[HEAL]}+other tag, +2 DV surcharge (minimum DV 4). Accept possible complications from SB. \\
Heal 2+ Harm (magic) & High DV (5+), often a ritual. Expect significant cost (Obligation, Corruption, quest). \\
Remove a Condition & \textit{Knot of Mercy} (Hedge Gift), \textit{Lay on Hands} (if Condition is curse/disease), or a Patron's Rite. \\
\bottomrule
\end{tabularx}
\end{table}

\section{Final Advice for Players}
\label{sec:healing-advice}

\begin{itemize}
  \item \textbf{Don"t hoard your healer's resources.} Fatigue is cheap to heal; use Mend often. Save big magic for when someone is at Harm 2 or unconscious.
  \item \textbf{Embrace the scars.} A permanent mark from a healing complication is a badge of survival. Work with your GM to make it meaningful.
  \item \textbf{Ambient healing is a thing.} After a fight, take a short rest. The GM will usually allow it unless you are under active pursuit.
  \item \textbf{Protect your healer.} If the only person with Mend goes down, your party's recovery slows to a crawl.
  \item \textbf{Talk about tone.} In Session Zero, discuss how gruesome or light you want healing complications to be. Some tables love body horror; others prefer comedic side‑effects.
\end{itemize}

\noindent The blade cuts. The spell mends. But the scar remains -- and that scar is your story.

